

%========================================================================
\section{Experimental Design}
\label{sec:model-setup}
%========================================================================
% One shared protocol for all models -> a fair comparison. This section
% describes HOW models are trained and evaluated; the architectures are in
% \cref{sec:models}, the conformal layer in \cref{sec:conformal}.

\subsection{Forecasting task}
\label{sec:exp:task}
For each country $i$ and decision month $t$, the models produce three
conditional quantiles of the one-month-ahead spread change $\dspread_{i,t+1}$
at levels $\mathcal{T} = \{0.1, 0.5, 0.9\}$, giving a central
$1-\alpha = 80\%$ prediction interval ($\alpha = 0.20$). The $0.5$ quantile
serves as the point estimator. The pair $\{0.1, 0.9\}$ was chosen because the
two pretrained models in the comparison, TiRex and TiRex-2, can only output the
deciles $\{0.1, 0.2, \dots, 0.9\}$; $0.1$ and $0.9$ are therefore the widest
levels available to every model without extrapolation. The trained quantile models, the deep models and LightGBM, are fitted with the
pinball (quantile) loss:
\begin{equation}
  \mathcal{L}(\hat{q}, y) \;=\;
  \frac{1}{|\mathcal{T}|}\sum_{\tau \in \mathcal{T}}
  \max\!\big(\tau\,(y - \hat{q}_\tau),\; (\tau - 1)(y - \hat{q}_\tau)\big).
  \label{eq:pinball}
\end{equation}
The point-forecast models (Random Walk and ARIMA) instead build their interval
by adding the empirical quantiles of their in-sample residuals to the point
forecast, which needs no distributional assumption.

\subsection{Training regimes}
\label{sec:exp:regimes}
The deep models are trained in three regimes:
\begin{itemize}
  \item \textbf{L} (local): one model per country, trained on up to 277 monthly
        observations per country.
  \item \textbf{G} (global): one model on the pooled panel of up to 2770
        country-months. The country identity is passed to the model, as a
        learned embedding for the neural networks and as a native categorical
        feature for LightGBM.
  \item \textbf{GF} (global fine-tuned): the global model is fully fine-tuned
        for each country at one tenth of the global training learning rate.
\end{itemize}
This way we get a fair comparison between a local, a global and a specialized model.
LightGBM uses only L and G, while the classical and zero-shot models use a
single setting (\cref{tab:model_selection}).

\subsection{Input sequences and scaling}
\label{sec:exp:sequences}
The input is the sequence of the last $W$ monthly feature vectors up to the
decision month $t$, from which we estimate $\dspread_{t+1}$. Before feeding it to
the model, we apply a per-country, per-feature (and per-target)
$z$-standardization, fitted only on the training window of each refit, so no
future information enters the scaling. The single missing input (Greece,
2015-07) is carried forward one month; the target and the evaluation are never
imputed.

\subsection{Rolling refit and temporal evaluation}
\label{sec:exp:refit}
We retrain all models on an expanding training set every 12 months. With the
training cut off at month $C$, the same model forecasts months
$C+1, C+2, \dots, C+12$ one step ahead, each using only the data available at
that point. The first 114 months form the minimal training set. Months 115 to
150 serve a double role: they are the burn-in set on which the conformal methods
are calibrated (\cref{sec:conformal}), and the validation set for the
hyperparameter optimization. Evaluation takes place on months 151 to 277, that
is 2015-10 to 2026-04, giving 127 test months per country.

\begin{table}[htbp]
    \centering
    \small
    \caption{Refit block structure. The cutoff month ends the training window, and the model then forecasts the listed months one step ahead.}
    \label{tab:refit_blocks}
    \begin{tabular}{@{}llll@{}}
        \toprule
        \textbf{Phase} & \textbf{Cutoff months} & \textbf{Forecast months} & \textbf{Evaluated} \\
        \midrule
        Burn-in & 114, 126, 138            & 115--150 & no  \\
        Test    & 150, 162, \dots, 270     & 151--277 & yes \\
        \bottomrule
    \end{tabular}
\end{table}

\subsection{Hyperparameter selection}
\label{sec:exp:hpo}
The hyperparameter optimization of all tuned models uses months 1 to 114 as the
training set and months 115 to 150 as the validation set. Afterwards the
hyperparameters are frozen and the models are retrained with them; no
test-window information enters the hyperparameter selection. Every trainable model with an information horizon is restricted to a shared maximum of 24 months, so that no model can use a longer lookback than the others. TiRex and TiRex-2 use the full expanding level history, which gives an information advantage to these models. Wherever a loss had to be minimized, the selection criterion was the mean pinball loss on the standardized $\Delta$ scale, averaged across countries, over the burn-in validation months.

\begin{table}[htbp]
  \centering\small
  \caption{Model selection across families. Info-horizon HP names the hyperparameter that sets a model's lookback; ``---'' marks a setting that does not exist for that family.}
  \label{tab:model_selection}
  \resizebox{\textwidth}{!}{%
  \begin{tabular}{@{}lllll@{}}
    \toprule
    \textbf{Family} & \textbf{Selection} & \textbf{Regimes} & \textbf{Info-horizon HP} & \textbf{Quantiles} \\
    \midrule
    Random Walk & none (ARMA(0,0) nesting)        & L        & ---                                   & empirical residuals \\
    ARMA        & AIC grid, $(p,q)\in\{0..3\}^2$   & L        & $(p,q)$ orders                        & empirical residuals \\
    LightGBM    & Optuna TPE, 100 trials           & L, G     & lag set $\subseteq\{1,2,3,6,12,24\}$  & quantile objective \\
    LSTM        & Optuna TPE, 100 trials           & L, G, GF & lookback $\in\{3,6,12,18,24\}$        & pinball loss \\
    xLSTM       & Optuna TPE, 100 trials           & L, G, GF & lookback $\in\{3,6,12,18,24\}$        & pinball loss \\
    TiRex, TiRex-2 & none (zero-shot, pretrained)  & ---      & native context window                 & native quantiles \\
    \bottomrule
  \end{tabular}%
  }
\end{table}

\paragraph{LSTM, xLSTM.}
Both share the same protocol and differ only in architecture. The
hyperparameters are tuned with Optuna (TPE sampler, seed 42, 20 start-up trials,
MedianPruner with 20 start-up trials and 3 warm-up steps), using 100 trials per
study and one study per regime (L, G). Within a trial, training runs for at most
100 epochs with early stopping (patience 15). The epoch count of the winning
trial is frozen (regime L: the median of the ten countries' best epochs; regime
G: the best epoch), and all refits train that same epoch count without early
stopping. The GF model is obtained by fine-tuning the G model at one tenth of the
learning rate; the number of fine-tuning epochs is chosen once on the burn-in
from the grid $\{5,10,20,40\}$. Three seeds (42, 43, 44) are run independently.
(Selected: xLSTM L $W{=}3$/5 ep., G $W{=}24$/1 ep.; LSTM L $W{=}3$/29 ep., G
$W{=}3$/10 ep.)

\paragraph{LightGBM.}
We use the same Optuna protocol (TPE, 100 trials, the same pruner, the same
burn-in selection), but only the regimes L and G exist: gradient boosting has no
fine-tuning analogue, so there is no GF. One booster is fitted per quantile.
Since LightGBM cannot process sequential data, we replace the lookback by an
ex-ante tabular lag structure: nested lag sets drawn from $\{1,2,3,6,12,24\}$
plus an optional block of rolling means over $\{3,6,12\}$ months, under the same
24-month cap. In the G regime the country enters as a native categorical
feature. Early stopping here acts on boosting rounds.

\paragraph{ARIMA, Random Walk.}
ARIMA has only the L regime and is therefore fitted per country on the
differenced target, which makes it an ARIMA$(p,1,q)$ in levels. For a given
order, the coefficients are estimated by maximum likelihood (exact, via the
Kalman filter); the order $(p,q)$ is selected by AIC over the grid
$(p,q)\in\{0,\dots,3\}^2$ (16 candidates) on the burn-in (data up to month 114)
and frozen per country. The models are therefore not fitted to a point-error
loss such as MAE or RMSE. The trend is suppressed (\texttt{trend="n"}), so
ARMA$(0,0)$ nests the Random Walk. Both models are deterministic, and their
prediction quantiles come from the empirical quantiles of their in-sample
one-step residuals, not from a quantile loss. They share the test window and the
conformal layer with the other models, but none of the selection machinery.

\paragraph{TiRex, TiRex-2.}
TiRex and TiRex-2 are pretrained time-series foundation models, so they need no
training and no hyperparameter tuning. Both are used exactly as released by NX-AI
on Hugging Face (\texttt{NX-AI/TiRex}, \texttt{NX-AI/TiRex-2}), and we use their
probabilistic output directly. TiRex-2 is additionally run in a covariate
variant; the variants are described in \cref{sec:models}.

\subsection{Seeds and aggregation}
\label{sec:exp:seeds}
Only the neural models, xLSTM and LSTM, use the three seeds (42, 43, 44). With
them we measure how much of the result is training randomness. Coverage is
reported as mean $\pm$ sd across seeds, and downstream we also report the span
of the median skill across seeds as a robustness check. The hyperparameters are
tuned once at seed 42 (\cref{sec:exp:hpo}) and then frozen, so the seeds only
differ in the production refits. The run loop iterates L, G, GF $\times$ blocks
$\times$ seeds. Every seed is its own independent refit/predict chain and keeps
its seed label on every prediction row through the conformal and diagnostic
stages. Predictions are never pooled across seeds before the per-chain metrics
are computed. The level-scale metrics (Winkler score, pinball loss, average
width, and median MAE) are computed per country, so per regime, seed, and
country, since otherwise a single high-spread country would dominate the
aggregate. Only the scale-free metrics, mean coverage rate and mean crossing
rate, are averaged across countries; this is the sole cross-country
aggregation. The remaining models, RW, ARIMA, LightGBM, TiRex, and TiRex-2,
are deterministic and run with a single seed.



\subsection{Quantile post-processing}
\label{sec:exp:rearrange}
The three quantiles are estimated independently, so they can cross. Crossed
quantiles would leave the conformal layer downstream with ill-defined interval
bounds. That is why we sort the three values per forecast before evaluation, a
monotone rearrangement \parencite{chernozhukov2010} that guarantees
$\hat{q}_{0.1} \le \hat{q}_{0.5} \le \hat{q}_{0.9}$. How often the raw
quantiles crossed before sorting is logged as a diagnostic and reported per
model. This step is the same for every model.

\subsection{Conformal prediction layer}
\label{sec:exp:conformal-ref}
The benchmarks emit only raw quantiles and contain no calibration code. The
conformal layer runs downstream on every combination of model, regime, seed, and
country. All methods share one nonconformity score, the CQR score
$E_t = \max(\hat{q}_{\mathrm{lo},t} - y_t,\; y_t - \hat{q}_{\mathrm{hi},t})$,
calibrated on a rolling out-of-sample pool. Since the same set of methods is
applied to every model, differences in interval quality can be traced back to
the models and not to the calibration. The methods, their fixed parameters, and
the calibration-pool variants are described in \cref{sec:conformal}.